\begin{topic}[id=sidechannel_attacks,
             difficulty={Anspruchsvoll},
             type={Sicherheitsanalyse + Bewertung},
             prerequisites={Grundlagen Kryptografie und Embedded Systems; Bereitschaft zur Einarbeitung in Mess- und Leckmodelle},
             practice={optional},
             owner={Dozententeam},
             updated={2026-04-09},
             tags={ SCA, DPA, CPA, EM, Timing, Gegenmaßnahmen, TVLA, Messaufbau }]{Seitenkanal-Angriffe auf Embedded-Systeme}

\TopicDescription{Seitenkanalangriffe werten unbeabsichtigte Signale aus—Stromaufnahme, EM-Abstrahlung oder Laufzeiten—und rekonstruieren daraus geheime Zustände. Eingebettete Systeme sind anfällig: wiederkehrende Muster, knappe Ressourcen und lange Lebenszyklen erschweren Härtung. Das Thema ordnet die wichtigsten Angriffsklassen (DPA/CPA, Timing, EM), erläutert Leck- und Auswertungsmodelle sowie typische Messaufbauten. Ebenso wichtig sind Gegenmaßnahmen: Maskierung, Randomisierung, konstante Zeit und physikalische Schutzmaßnahmen. Du erhältst einen nüchternen Blick auf Evaluierung (z.~B. TVLA), Kosten/Nutzen-Abwägungen und Designprinzipien für mehr Resilienz.}

\TopicLeitfragen{\item Welche Leckmodelle und Messgrößen sind praxistauglich?
\item Welche Gegenmaßnahmen sind wirksam bei vertretbarem Overhead?
\item Wie liest man TVLA-Ergebnisse richtig?
}
\TopicScope{\item Kein HF-Messtechnik-Deep-Dive.
\item Keine Kryptografie-Theorie im Detail.
}
\TopicPractice{Mini-Experiment (real/simuliert) zu DPA/CPA mit kurzer Auswertungsdemonstration.}

\TopicKeywords{SCA, DPA, CPA, EM, Timing, Gegenmaßnahmen, TVLA, Messaufbau}

\nocite{kocherDPA1999,chipwhispererDocs}
\end{topic}
